\section{Conclusão e Trabalhos Futuros}\label{sec:conc}
\AtBeginSec

%\begin{frame}{O que foi discutido}
%\begin{bigitem}
% \item Formas de aumentar a taxa de transmissão de dados. Ex: alocação de potência e emparelhamento de subportadoras.
% \item Problema conjunto de emparelhamento de subportadoras e alocação de potência para sistemas multicanais.
% \item O problema formulado é classificado como misto binário-inteiro e NP-difícil.
% \item Separamos em dois subproblemas independentes: emparelhamento de subportadoras e alocação de potência.
%\end{bigitem}
%\end{frame}
%
%\begin{frame}{O que foi solucionado}
%\begin{bigitem}
% \item Emparelhamento de subportadoras:
% \begin{bigitem}
% \item Qualquer emparelhamento será tão bom quanto ordenar as subportadoras dos dois saltos em ordem decrescente de seus ganhos de canal e então emparelhá-las.
% \end{bigitem}
% \item Alocação de potência:
% \begin{bigitem}
%  \item Relação linear entre a potência dos saltos, reduzindo a complexidade de sua solução.
% \end{bigitem}
% \item Após o tratamento independente, é possível unificá-los de forma sequencial\\
% emparelhamento de subportadoras $\Rightarrow$ alocação de potência
%\end{bigitem}
%\end{frame}
%
%\begin{frame}{O que foi desenvolvido}
% \begin{bigitem}
%  \item Duas relaxações do problema de alocação de potência, permitindo desenvolver dois algoritmos subótimos baseados em Water-Filling:
%   \begin{bigitem}
%    \item EOS+EP (Emparelhamento Ótimo de Subportadoras e Escalonamento de Potência).
%    \item EOS+APM (Emparelhamento Ótimo de Subportadoras e Alocação de Potência Mínima).
%   \end{bigitem}
% \end{bigitem}
%\end{frame}
%
%\begin{frame}{O que foi obtido}
%\begin{bigitem}
%  \item EOS+EP alcançou os melhores resultados dentre as estratégias subótimas em todos os cenários avaliados.
%  \item EOS+APM alcançou resultados próximos à solução ótima computacional mas pior que EOS+EP.
%  \item EOS+EP é uma ótima escolha quando o poder de processamento não é tão alto e o tempo de execução é relevante.
%  \item Há uma melhoria considerável no desempenho dos algoritmos ao utilizar EOS.
%\end{bigitem}
%\end{frame}

\begin{frame}{Conclusão}
\begin{bigitem}
  \item Problema conjunto de emparelhamento de subportadoras e alocação de potência para sistemas multicanais.
  \item Separamos em dois subproblemas independentes: emparelhamento de subportadoras e alocação de potência.
  \item Após o tratamento independente, é possível unificá-los de forma sequencial\\
 emparelhamento de subportadoras $\Rightarrow$ alocação de potência
  \item Duas relaxações do problema de alocação de potência, permitindo desenvolver dois algoritmos subótimos baseados em Water-Filling.
  \item Os algoritmos alcançaram ótimos resultados em todos os cenários analisados.
\end{bigitem}
\end{frame}

\begin{frame}{Perspectivas futuras}
\begin{bigitem}
  \item Sistemas de transmissão e recepção MIMO com conhecimento perfeito da CSI.
  \item Esperamos poder dividir o problema em dois independentes e após propor uma solução, unificá-los novamente.
\end{bigitem}
\end{frame}

\begin{frame}
  \begin{figure}[!htb]
    \centering
    \includegraphics[width=0.60\linewidth]{../Imagens/obrigado.eps}    
   \end{figure}
\end{frame}